\section{External gear-rig comparisons}

The simulated gear benchmarks above control numerical resolution and the contact representation. We additionally test the reduced contact dynamics against a single-tooth force measurement and against two-shaft encoder measurements. Stiffness and damping are selected from a fixed grid using one operating condition per rig. The mesh and B-rep runs share the selected parameters. Scoring on held-out conditions uses the same parameters and phase conventions.

\subsection{NASA single-tooth dynamic force}

We use the reconstructed 28/28-tooth CAD fixture described above and constant torque at the values indicated in Fig.~7 of NASA TM-103232. Four torsional coordinates represent the motor, pinion, wheel, and load. The model uses the inertia values printed in SI units in Table II and the reported 341 and 12,700 N\,m/rad input and output torsional stiffnesses. An equivalent normal compliance represents the working contact. The gearbox stiffness entry from the reference is distinct from the effective normal tooth stiffness fitted here. The shaft damping is set to 0.05 of critical for the adjacent-inertia pair. The selected effective tooth stiffness is \NasaK{} MN/m with damping ratio \NasaZeta{}. Both values lie on the prescribed NASA calibration grid boundary.

We digitize the four experimental traces, using only the solid force curves. The points trace the published plot with an ordinate resolution of roughly 6 N per source pixel. Thick and overlapping strokes add local reading uncertainty. The comparison uses the driven tooth roll angle from the actual native or meshed witness. We use the simulated single-pair forces, average complete simulated engagements in the last half of each run, and hold the roll-angle origin fixed. The scored range ends at \(29.5^\circ\), with the lower bound at \(12^\circ\) for (a)--(c) and \(12.625^\circ\) for (d). The restricted endpoint of (d) follows the identifiable trace of the plotted measurement. \href{https://ntrs.nasa.gov/citations/19910013253}{NASA TM-103232}.

\begin{table}[!htbp]
\centering\small
\caption{Single-tooth normal force relative to the digitized experiments. Peak refers to the scored roll-angle range. RMSE is computed on the common \(0.125^\circ\) grid.}
\label{tab:nasa-exp}
\begin{tabular}{lrrrrrl}
\toprule
Panel & rpm & Torque (N\,m) & Measured peak (N) & B-rep RMSE (N) & Mesh RMSE (N) & Role\\
\midrule
\NasaExperimentRows
\bottomrule
\end{tabular}
\end{table}

\begin{figure}[!htbp]
\centering
\includegraphics[width=.88\linewidth]{\GearExperimentPath nasa-single-tooth.pdf}
\caption{NASA single-tooth force. Dots digitize the experimental trace. Solid and dashed lines are the HHT runs using native STEP faces and triangles, respectively.}
\label{fig:nasa-exp}
\end{figure}

The low-load 4000 rpm measurement shows loss of tooth contact. Within the central \(19.5^\circ\) to \(21.5^\circ\) interval the native single-tooth force remains above \NasaLowLoadFloor{} N, leaving the loss of contact unpredicted. In the 6000 rpm, 94\% torque case, halving the native time step changes the force profile by \NasaStepDelta{} N RMS. The paired SI and customary inertia entries in NASA Table II imply a tenfold difference when the customary entries are interpreted as force-inch units. A separate sensitivity run multiplies all four inertias by ten with the other settings fixed. At 6000 rpm and 47\% torque it changes the force profile by \NasaInertiaSensitivity{} N RMS, giving an error of \NasaInertiaRMSE{} N against the digitized measurement. The main comparison retains the published SI entries.

\IfFileExists{nasa-revised-results.tex}{%
  % Generated by benchmarks/gears/experimental/nasa_revision_artifacts.py.
\newcommand{\NasaRevisedPath}{\detokenize{../results/gears/experimental/revised/static-matched/figures/}}
\newcommand{\NasaRevisedStiffness}{67.5}
\newcommand{\NasaRevisedDamping}{0.10}
\newcommand{\NasaRevisedKBeam}{172.8}
\newcommand{\NasaRevisedFrequencies}{6.73, 665.63, 4106.59}
\newcommand{\NasaRevisedGearboxFrequencies}{6.56, 524.90, 1220.79}
\newcommand{\NasaRevisedStaticSpan}{15.87, 24.88}
\newcommand{\NasaRevisedStaticTarget}{16.00, 25.00}
\newcommand{\NasaRevisedLowMin}{125.9}
\newcommand{\NasaRevisedLowOverlap}{0}
\newcommand{\NasaRevisedObservedSpan}{19.25, 22.00}
\newcommand{\NasaBeamInternalSpan}{23.50, 24.38}
\newcommand{\NasaRevisedSmallestGap}{29.33}
\newcommand{\NasaRevisedStepB}{0.42}
\newcommand{\NasaRevisedStepD}{0.16}
\newcommand{\NasaQuarterStepDelta}{1.33}
\newcommand{\NasaQuarterStepMin}{125.21}
\newcommand{\NasaQuarterStepSamples}{51,409}
\newcommand{\NasaRevisedTimeBrep}{1.53}
\newcommand{\NasaRevisedTimeBvh}{6.94}
\newcommand{\NasaRevisedSpeed}{4.54}
\newcommand{\NasaRevisedDualShares}{10.3, 29.9, 60.7}
\newcommand{\NasaRevisedRows}{%
7(a) & 63.1 & 67.8 & 125.4 & 63.2 & 63.9 & -3.5 \\
7(b) & 146.7 & 173.7 & 226.4 & 203.2 & 201.9 & +0.6 \\
7(c) & 148.8 & 164.7 & 332.6 & 174.0 & 176.0 & -17.6 \\
7(d) & 90.8 & 90.7 & 260.1 & 75.7 & 81.6 & -4.7 \\
}
%
  \subsection{Reconstructed gear compliance and the loss-of-contact test}

We rebuilt both STEP flanks from the average shape of the four traces per gear in Fig.~3 of NASA TM-103232. The plotted profiles provide the relative shape. Table~I fixes the tip amplitudes at 23 and \(25\,\mu\mathrm m\). This reconstruction shares the same geometric profile across all teeth of a gear. The native and triangle runs use the same exported STEP surfaces and trimmed line-of-action stations. We converted the customary inertia column to SI using \(1\,\mathrm{lbf\,in\,s^2}=0.112985\,\mathrm{kg\,m^2}\). The \href{https://ntrs.nasa.gov/citations/19930018861}{DANST manual} independently lists the customary magnitudes. The equivalent gearbox stiffness in Table~II includes shaft flexibility. It enters the modal audit separately from the tooth pair compliance.

At each station we compute the compliance of the two gear teeth as tapered beams with bending, shear, and radial compression, and add a line-contact compliance. Let \(g_i\) denote the CAD gap, \(g_0\) the fixed gap offset, \(C_i\) the sum of those compliances, and \(C_b\) a constant series compliance accounting for deformation outside the reduced tooth model. With geometric endpoint weight \(w_i\), rolling normal force satisfies
\[
 k_i=(C_i+C_b)^{-1},\qquad
 F_i=w_i\max\!\left(0,k_i(g_0-g_i)-c_i\,\dot g_i\right)
 \quad\text{for }g_i<g_0 ,
\]
and vanishes for \(g_i\ge g_0\). The damping coefficient uses a ratio of \NasaRevisedDamping{} and the reduced inertia. The gap derivative includes the advance of the nominal gear angle and the generalized velocities. We obtain its signed lever arms from the tangent of each native isocurve or the normal of the intersected triangle. The load tangent for the HHT solve uses the same gap gradient. The query keeps the finite-flank trim and the original \(2\,\mu\mathrm m\) sampled working-flank guard.

The independent beam and line-contact estimate gives \NasaRevisedKBeam{} MN/m at the pitch point. For the series-compliance variant we solve torque balance over 224 mesh phases at each load. We select one stiffness using only the approximate static DANST plateau in Fig.~7(c), near roll angles \(\NasaRevisedStaticTarget\) degrees. The selected pitch stiffness is \NasaRevisedStiffness{} MN/m, with plateau endpoints \(\NasaRevisedStaticSpan\) degrees. The reference plateau is a published calculation. Its overlapping strokes limit the precision of the endpoints. The static torque balance residual stays below \(10^{-6}\,\mathrm{N\,m}\). For the native runs, the fractions of phases with two loaded tooth pairs are \NasaRevisedDualShares{} percent at 11.48, 33.73, and 67.46 N\,m, respectively.

The four dynamic comparisons use the same stiffness and damping for both geometric backends. The gearbox-stiffness linearization with the converted inertias gives \(\NasaRevisedGearboxFrequencies\) Hz. The report lists 6.56, 52.5, and 1220 Hz. The middle entry is inconsistent with our four-inertia reconstruction. The single loaded pair at the calibrated pitch stiffness gives \(\NasaRevisedFrequencies\) Hz. Table~II lists the highest DANST frequency at 5821 Hz. A sensitivity run imposed that value on the frozen single-pair linearization. It produced an early zero-force interval with zero overlap with the measured zero bins.

\begin{figure}[!htbp]
\centering
\includegraphics[width=.87\linewidth]{\NasaRevisedPath nasa-revised-force.pdf}
\caption{NASA single-tooth force with converted inertias, digitized profile shapes, and the static-calibrated series compliance. The shaded region in 7(b) marks the measured zero-force interval on the scoring grid.}
\label{fig:nasa-revised}
\end{figure}

\begin{table}[!htbp]
\centering\small
\caption{Waveform RMSE in newtons for the successive NASA model changes. Peak error refers to the static-calibrated native model.}
\label{tab:nasa-revised}
\resizebox{\linewidth}{!}{%
\begin{tabular}{lrrrrrr}
\toprule
Panel & Original B-rep & Inertias only & Beam B-rep & Calibrated B-rep & Calibrated BVH & B-rep peak err. (\%)\\
\midrule
\NasaRevisedRows
\bottomrule
\end{tabular}}
\end{table}

In 7(b) the digitized measured zero range is \(\NasaRevisedObservedSpan\) degrees. The pure beam model predicts an internal zero-force interval at \(\NasaBeamInternalSpan\) degrees. The calibrated B-rep result overlaps \NasaRevisedLowOverlap{} percent of the measured zero bins and reaches a minimum of \NasaRevisedLowMin{} N on them. The observed phase discrepancy persists under step refinement. Halving the step changes the calibrated B-rep waveform by \NasaRevisedStepB{} N RMS in 7(b) and \NasaRevisedStepD{} N RMS in 7(d). The calibrated runs retain a smallest recorded gap of \NasaRevisedSmallestGap{} micrometers. The remaining force errors motivate an independently identified dynamic model of the shafts and supports.

Refining 7(b) again to \(0.25\,\mu\mathrm s\) changes its waveform by \NasaQuarterStepDelta{} N RMS relative to the nominal step. Before phase averaging, \NasaQuarterStepSamples{} saved tooth samples occupy the measured zero-force interval. Their minimum force is \NasaQuarterStepMin{} N. The output cadence in this audit is \(1\,\mu\mathrm s\).

In three alternating, serial timing repeats of 7(b), each simulating 0.16 s, median HHT integration takes \NasaRevisedTimeBrep{} s for the native surfaces and \NasaRevisedTimeBvh{} s for the triangle BVH, a factor of \NasaRevisedSpeed{}. Timings exclude geometry setup and output.
%
}{}

\subsection{UNSW encoder transmission error}

The UNSW test gear pair has 27/44 teeth, a 2 mm module, 5 mm face width, \(20^\circ\) pressure angle, and 10 mm bore radii. We use three healthy records from the original 200 kHz dataset. The input reference phase is set by the tachometer. The 3600-pulse encoder signals on both shafts are converted to continuously interpolated angular measurements using linearly interpolated zero crossings. The base-radius-scaled difference gives the transmission error. We remove the block mean and average complete hunting-tooth periods of 44 input revolutions, preserving the mesh-synchronous waveform. Repeating the extraction with falling edges changes the three mean waveforms by less than \(0.012\,\mu\)m RMS. The reported block range shows variation within each recording.

The reconstructed solver uses two effective shaft inertias and ideal involute STEP geometry. Encoder transfer dynamics and as-manufactured deviations are omitted from this model. The effective stiffness \UnswK{} MN/m, damping ratio \UnswZeta{}, and the phase offset of \UnswShift{} mesh cycles are set using the nominal 5 Hz and 10 N\,m record. The two held-out records use measured speeds and mean input torques. Halving the native step in the faster held-out case changes the mesh-synchronous waveform by \UnswStepDelta{} \(\mu\)m RMS. \href{https://doi.org/10.26190/UNSWORKS/24007}{Test geometry in Chin's thesis}. \href{https://doi.org/10.17632/p72x3m92cv.1}{CC BY 4.0 dataset}.

\begin{table}[!htbp]
\centering\small
\caption{UNSW AC transmission error. RMS and RMSE are computed on 128 uniformly spaced points over one mesh cycle.}
\label{tab:unsw-exp}
\begin{tabular}{rrrrrl}
\toprule
Measured rpm & Input N\,m & Data RMS (\(\mu\)m) & B-rep RMSE (\(\mu\)m) &
Mesh RMSE (\(\mu\)m) & Role\\
\midrule
\UnswExperimentRows
\bottomrule
\end{tabular}
\end{table}

\begin{figure}[!htbp]
\centering
\includegraphics[width=.88\linewidth]{\GearExperimentPath unsw-encoder-dte.pdf}
\caption{UNSW DTE and mesh harmonics at the calibration setting and two held-out settings. The shaded band is the range across complete hunting-tooth blocks of the encoder record.}
\label{fig:unsw-exp}
\end{figure}

At the faster held-out setting the measured second mesh harmonic is \UnswMeasuredSecond{} \(\mu\)m and the native result is \UnswNativeSecond{} \(\mu\)m. The held-out errors characterize the full reduced model. Residual discrepancies can include contact-law approximations, rig dynamics, and manufacturing deviations. Replacing the triangulated faces with native STEP geometry isolates one numerical contribution within those settings.
